\documentclass{article}
\usepackage{amsmath}

\title{Ch. 5.2-5.3 HW Solns}
\author{Mr. Spence}
\date{\today}

\begin{document}

\maketitle

\subsection*{QuickQuiz}

\begin{enumerate}
\setcounter{enumi}{4}

\item (b)

	\begin{gather*}
	\% \Delta P = \frac{24 - 16}{\frac{24 + 16}{2}} = \frac{8}{20} = \frac{2}{5} \\
	\% \Delta Q = \frac{110 - 90}{\frac{110 + 90}{2}} =  \frac{20}{100} = \frac{1}{5} \\
	\text{Price Elasticity of Supply} = \frac{\% \Delta Q}{\% \Delta P} = \frac{1}{2} = 0.5.
	\end{gather*}
	
\item (c) -- If the price elasticity of supply is zero, that means the supply curve is perfectly inelastic, and perfectly inelastic supply is vertical.

\item (c) -- The ability of firms to enter and exit a market in the long run gives sellers more flexibility to change the amount they produce, so it makes the supply curve more elastic.

\end{enumerate}

\subsection*{Problems and Applications}

\begin{enumerate}
\setcounter{enumi}{7}

\item

	\begin{enumerate}
	
	\item 
	
		\begin{gather*}
		\% \Delta Q = -0.043 \\
		\% \Delta P = \frac{1.5 - 1.25}{\frac{1.5 + 1.25}{2}} = \frac{0.25}{1.375} \approx 0.18 \\
		\eta = \left| \frac{\% \Delta Q}{\% \Delta P} \right| \approx 0.24
		\end{gather*}
		
	\item Based on my estimate, demand for subway rides is inelastic, so when the Transit Authority raises fares, we should expect revenue to increase. That's because quantity demanded will decrease proportionately less than fares will increase. 
	
	\item The estimate might be unreliable because there could be other reasons for the decrease in ridership in December. For example, it's possible that so many people take vacation in December that subway ridership drops significantly. (This question doesn't test anything we've learned. It's just a common sense question that hints at econometrics, a statistics course you'll take in college if you choose to major in economics.) 
	
	\end{enumerate} 
	
\item Thelma's demand is perfectly inelastic ($\eta = 0$) because she demands the same quantity regardless of price. Louise's demand is perfectly elastic ($\eta = \infty$) because she's willing to buy any quantity at \$20. 

\item

	\begin{enumerate}
	
	\item 
	
		\begin{gather*}
		\eta = \left| \frac{\% \Delta Q}{\% \Delta P} \right| \\
		0.4 = \left| \frac{-20 \%}{\% \Delta P} \right| \\
		\% \Delta P = \frac{20}{0.4} \\
		\% \Delta P = 50 \\
		100 * \frac{P_2 - 5}{\frac{P_2 + 5}{2}} = 50 \\
		P_2 - 5 = 0.25P_2 + 1.25 \\
		0.75P_2 = 6.25 \\
		P_2 \approx \$8.33
		\end{gather*}
		
		The government should increase the price by \$3.33.
		
	\item The policy will have a larger effect five years from now. In the long run, demand will be more elastic, so quantity demanded will decrease more in response to the price increase. 
	
	\item Teenager smokers might have developed less of a nicotine addiction than adult smokers. As a result, they might treat cigarettes as more of a luxury good while adults might treat them as more of a necessity. (I didn't mention this in class, so you don't need to know it, but generally, when spending on a good constitutes a high share of people's income, demand for that good is more elastic. Teenagers tend to have low incomes, so cigarettes constitute a higher share of their income.)
	
	\end{enumerate}


\end{enumerate}

\end{document}